\mode*
\subsection{Files}
\label{files}

In our course, files come late, after strings, containers and error
handling, and a file program reuses all three: a file is opened, read into strings, and its
failures arrive as exceptions.
The literature on novice misconceptions says little about files as such.
A survey of three decades of publications on the logic errors of CS1
students sorts the errors it found into eleven categories ---
\textcquote{AlzahraniVahid2021}{input (2 errors), output (1 error),
variable (7 errors), computation (21 errors), condition (18 errors), branch
(14 errors), loop (27 errors), array (5 errors), function (24 errors),
conceptual (43 errors), and miscellaneous (4 errors)} --- and none of them
concerns files.
Our own two-sided search found no study of novices' conceptions of opening,
reading or writing files either.
The analyses below therefore rest on two kinds of ground: misconceptions
documented for the constructs a file program reuses (exceptions, and
values read as strings), and, where we found no study, the critical
aspects that our course material on files is designed around.
We say which is which in each case.

\paragraph{Closing a file when something goes wrong.}
A file opened with \mintinline{python}{open} must be closed again, and
the simplest program does so on the line after it has used the file.
If anything between the two lines raises an exception, the closing line
never runs.
Seeing this requires what \textcite{RashkovitsLavy2012} found missing in
students at the end of a Java programme: they
\textcquote[343]{RashkovitsLavy2012}{had not assimilated the fact that once an
exception is thrown, the program jumps to the closest catch clause that is appropriate for the exception and continues from there}, and
\enquote{flow of control in the context of exceptions} was among the
aspects they found hardest \autocite[350]{RashkovitsLavy2012}.
The study concerns Java and students near graduation, not first-year
Python students; that the mechanism is the same in both languages is our
inference.

The object of learning is \emph{why a file opened by hand may stay open}.
The critical aspect is that the exception skips the rest of the protected
block, including the \mintinline{python}{close}; the
\mintinline{python}{with} statement closes the file however its block is
left \autocite{PythonTutorialFiles}.
We hold the file (it contains the word \texttt{forty-two}, so
\mintinline{python}{int} fails), the handler and the final question
invariant, and vary only how the file is closed.

\begin{description}
  \item[Contrast] The left program prints \mintinline{python}{False}: the
    \mintinline{python}{close} after the failing \mintinline{python}{int}
    never ran.
    The right one prints \mintinline{python}{True}.
    Since everything else is the same, only the way of closing can explain
    the difference --- and it shows only when something goes wrong.

    \hfill
    \begin{minipage}[t]{0.45\columnwidth}
      \begin{minted}[highlightlines={1,4}]{python}
        f = open("data.txt")
        try:
            print(int(f.read()))
            f.close()
        except ValueError:
            print("Not a number")
        print(f.closed)
      \end{minted}
    \end{minipage}
    \hfill
    \begin{minipage}[t]{0.45\columnwidth}
      \begin{minted}[highlightlines={2}]{python}
        try:
            with open("data.txt") as f:
                print(int(f.read()))
        except ValueError:
            print("Not a number")
        print(f.closed)
      \end{minted}
    \end{minipage}
\end{description}

\paragraph{A file that does not exist.}
Opening a file that does not exist raises
\mintinline{python}{FileNotFoundError}, the exception Python raises
\textcquote{PythonExceptions}{when a file or directory is requested but
doesn't exist}, and it is \mintinline{python}{open} that raises it, before
anything is read.
Both parts carry a belief we found no study of: that a missing file reads
as an empty one, as if \mintinline{python}{open} for reading created it the
way opening for writing does.
The useful response is to catch the exception where the program knows what
to do instead --- typically to ask for another name --- which is usually not
the function that opens the file.
That step needs another of the aspects \textcite{RashkovitsLavy2012} found
hard, \textcquote[350]{RashkovitsLavy2012}{handling exceptions further up the
calling chain}.

The object of learning is \emph{what happens when the file is missing}.
The critical aspects are that the failure is an exception raised by
\mintinline{python}{open}, not an empty file, and that it travels out of
the reading function to whoever catches it.
We hold the program invariant and vary only whether the file exists.

\begin{description}
  \item[Contrast] The same program, run once in a folder where
    \texttt{notes.txt} exists and once where it does not.
    With the file, it prints \texttt{Got the text}; without it,
    \mintinline{python}{open} raises inside \mintinline{python}{read_text},
    the rest of the \mintinline{python}{try} block is skipped, and the
    caller's handler prints \texttt{No such file}.
    Both runs end with \texttt{Done}.

    \hfill
    \begin{minipage}[t]{0.6\columnwidth}
      \begin{minted}[highlightlines={2,8}]{python}
        def read_text(filename):
            with open(filename) as f:
                return f.read()

        try:
            text = read_text("notes.txt")
            print("Got the text")
        except FileNotFoundError:
            print("No such file")
        print("Done")
      \end{minted}
    \end{minipage}
    \hfill\null
\end{description}

\paragraph{The line keeps its line break.}
Looping over a file gives its lines, and each line ends with the
line-break character that separated it from the next: Python's tutorial
notes that \textcquote{PythonTutorialFiles}{a newline character (\textbackslash
n) is left at the end of the string}.
A student who expects the line without it will find that
\mintinline{python}{print(line)} double-spaces the output and that
\mintinline{python}{line == "Pippi"} is never true.
We found no study documenting this belief; it is the aspect our course
material targets when it asks students to write the loop themselves before
showing its output.

The object of learning is \emph{what one step of the loop over a file
gives}.
The critical aspect is that the unit is the line \emph{including} its
\mintinline{python}{"\n"} --- not the line without it, not the whole file,
and not a single character as when looping over a string.
We hold the file and the loop invariant and vary only whether
\mintinline{python}{print} adds a line break of its own.

\begin{description}
  \item[Contrast] The file holds the two lines \texttt{tea} and
    \texttt{jam}.
    The left loop prints each line followed by an empty line, the right one
    prints them as they stand in the file.
    The only difference is \mintinline{python}{end=""}, so the empty lines
    can only come from the line breaks the lines carry.

    \hfill
    \begin{minipage}[t]{0.45\columnwidth}
      \begin{minted}[highlightlines={3}]{python}
        with open("food.txt") as f:
            for line in f:
                print(line)
      \end{minted}
    \end{minipage}
    \hfill
    \begin{minipage}[t]{0.45\columnwidth}
      \begin{minted}[highlightlines={3}]{python}
        with open("food.txt") as f:
            for line in f:
                print(line, end="")
      \end{minted}
    \end{minipage}
\end{description}

\paragraph{Writing replaces, appending adds.}
Opening a file for writing empties it first: in the tutorial's words, with
mode \mintinline{python}{"w"} \textcquote{PythonTutorialFiles}{an existing
file with the same name will be erased}, whereas \mintinline{python}{"a"}
adds to the end.
A program that opens the same file for writing twice --- in a function
called twice, or inside a loop --- therefore keeps only what the last call
wrote.
Again we found no study of how novices conceive of the modes; that the
destruction is silent is what makes it worth a dedicated contrast.

The object of learning is \emph{what opening a file does to what the file
already holds}.
The critical aspect is the mode: \mintinline{python}{"w"} discards the old
content, \mintinline{python}{"a"} keeps it.
We hold the writer, the two calls and the reading invariant and vary only
the mode.

\begin{description}
  \item[Contrast] Starting without a \texttt{log.txt}, after the two calls
    the left file holds one line, \texttt{two}, and the right one holds
    \texttt{one} and \texttt{two}.

    \hfill
    \begin{minipage}[t]{0.45\columnwidth}
      \begin{minted}[highlightlines={2}]{python}
        def save(text):
            with open("log.txt", "w") as f:
                f.write(text + "\n")

        save("one")
        save("two")
      \end{minted}
    \end{minipage}
    \hfill
    \begin{minipage}[t]{0.45\columnwidth}
      \begin{minted}[highlightlines={2}]{python}
        def save(text):
            with open("log.txt", "a") as f:
                f.write(text + "\n")

        save("one")
        save("two")
      \end{minted}
    \end{minipage}
\end{description}

\paragraph{What is read is text.}
The data-types analysis showed that \mintinline{python}{input} returns a
string even when the user types a number.
The same holds for files, and the misconception is documented for both
sources: \textcquote{GuoMarkelZhang2020}{When numerical data is read from
files or terminal input, they often start as strings}, and
\textcite{Bosse2021Antipatterns} catalogue the missing conversion as an
antipattern in CS1 code.
For files the error is easier to make, because the content looks like a
number and nobody typed it.

The object of learning, as for \mintinline{python}{input}, is \emph{that a
value's type is independent of how the value looks}.
The critical aspect is that reading a file gives strings, so the program
must convert what it reads.
We hold the file (it contains \texttt{42}) and the operation invariant and
vary only whether \mintinline{python}{int} is applied --- the
generalisation, from the terminal to the file, of the contrast in the
data-types analysis.

\begin{description}
  \item[Generalisation] The left program prints \texttt{4242}, the right
    one \texttt{84}.
    The source of the string has changed from the keyboard to a file; the
    critical aspect, and the conversion it requires, has not.

    \hfill
    \begin{minipage}[t]{0.45\columnwidth}
      \begin{minted}[highlightlines={2}]{python}
        with open("number.txt") as f:
            value = f.read().strip()
        print(value * 2)
      \end{minted}
    \end{minipage}
    \hfill
    \begin{minipage}[t]{0.45\columnwidth}
      \begin{minted}[highlightlines={2}]{python}
        with open("number.txt") as f:
            value = int(f.read().strip())
        print(value * 2)
      \end{minted}
    \end{minipage}
\end{description}
